\section{量子ゲートと複数量子ビット}
\label{sec:gates-composite}

\subsection{ユニタリ変換と1量子ビットゲート}

量子状態に行列を掛けると，別の量子状態へ変化する．
ただし，変化の前後でベクトルの長さが変わってはならない．
この条件を満たす行列または線形写像をユニタリという．

\begin{defi}[ユニタリ写像]
内積空間$\mathcal H$上の線形写像$U$が
$U^\dagger U=UU^\dagger=I$を満たすとき，$U$をユニタリ写像という．
\end{defi}

量子系が測定されず，外部から影響を受けない間，
その状態変化はユニタリ写像で表されるとする．
したがって，初期状態が$\ket{\psi}$なら，変化後の状態は$U\ket{\psi}$である．

\begin{prop}
ユニタリ写像$U$は内積とノルムを保存し，逆写像は$U^{-1}=U^\dagger$である．
\end{prop}
\begin{proof}
任意の$\ket{\phi},\ket{\psi}$について
\[
  \braket{U\phi|U\psi}
  =\bra{\phi}U^\dagger U\ket{\psi}
  =\braket{\phi|\psi}
\]
である．両ベクトルを$\ket{\phi}=\ket{\psi}$とすればノルムの保存が従う．
また，$U^\dagger U=UU^\dagger=I$である．
したがって，$U^\dagger$は$U$の逆写像である．
\end{proof}

HHLまでに用いる基本的な1量子ビットゲートを
\begin{align}
  X&=\begin{pmatrix}0&1\\1&0\end{pmatrix},&
  Y&=\begin{pmatrix}0&-\mathrm{i}\\\mathrm{i}&0\end{pmatrix},&
  Z&=\begin{pmatrix}1&0\\0&-1\end{pmatrix},\\
  H&=\frac1{\sqrt2}\begin{pmatrix}1&1\\1&-1\end{pmatrix}
\end{align}
で定める．直接の行列計算により
$X^\dagger X=Y^\dagger Y=Z^\dagger Z=H^\dagger H=I$であるから，
これらはすべてユニタリ写像である．特に
\begin{equation}
  H\ket{0}=\frac{\ket{0}+\ket{1}}{\sqrt2},
  \qquad
  H\ket{1}=\frac{\ket{0}-\ket{1}}{\sqrt2}
  \label{eq:hadamard-action}
\end{equation}
である．

等式$Y^2=I$を用いると，$y$軸回りの回転は
\begin{align}
  R_y(\theta)
  &:=\exp\left(-\frac{\mathrm{i}\theta}{2}Y\right)\\
  &=\cos\frac{\theta}{2}I-\mathrm{i}\sin\frac{\theta}{2}Y
   =\begin{pmatrix}
      \cos(\theta/2)&-\sin(\theta/2)\\
      \sin(\theta/2)& \cos(\theta/2)
    \end{pmatrix}
  \label{eq:ry-definition}
\end{align}
となる．第2の等号は指数関数の級数を偶数次と奇数次に分けて得られる．
式\eqref{eq:ry-definition}の実行列表示において$c=\cos(\theta/2)$，
$s=\sin(\theta/2)$とおけば，
\[
 R_y(\theta)^\dagger R_y(\theta)
 =\begin{pmatrix}c&s\\-s&c\end{pmatrix}
  \begin{pmatrix}c&-s\\s&c\end{pmatrix}
 =\begin{pmatrix}c^2+s^2&0\\0&c^2+s^2\end{pmatrix}=I
\]
であるから，$R_y(\theta)$もユニタリ写像である．
したがって
\begin{equation}
  R_y(\theta)\ket{0}
  =\cos\frac{\theta}{2}\ket{0}
   +\sin\frac{\theta}{2}\ket{1}
  \label{eq:ry-on-zero}
\end{equation}
である．

\subsection{テンソル積と複合系}

二つの量子ビットを同時に扱うには，二つのベクトル空間を組み合わせる必要がある．
この組合せに使う演算がテンソル積である．
まず，資料\cite[73--75頁]{tsukano2023}に従い，
行列のテンソル積をブロック行列として定義する．

\begin{defi}[行列のテンソル積]
$m\times n$行列$A=(a_{ij})$と$r\times s$行列$B$に対し，
\begin{equation}
 A\otimes B
 :=
 \begin{pmatrix}
  a_{11}B & a_{12}B & \cdots & a_{1n}B\\
  a_{21}B & a_{22}B & \cdots & a_{2n}B\\
  \vdots  & \vdots  & \ddots & \vdots\\
  a_{m1}B & a_{m2}B & \cdots & a_{mn}B
 \end{pmatrix}
 \label{eq:matrix-tensor-product-definition}
\end{equation}
を$A$と$B$のテンソル積という．これは$mr\times ns$行列である．
\end{defi}

列ベクトル$v=(v_1,\ldots,v_m)^{\mathsf T}$と
$w=(w_1,\ldots,w_r)^{\mathsf T}$も行列とみなせるため，
式\eqref{eq:matrix-tensor-product-definition}から
\begin{equation}
 v\otimes w
 =
 \begin{pmatrix}
  v_1w\\v_2w\\\vdots\\v_mw
 \end{pmatrix}
 \in\mathbb C^{mr}
 \label{eq:vector-tensor-product-coordinate}
\end{equation}
となる．つまり，各成分$v_i$をベクトル$w$全体に掛け，
その結果を上から順に並べる．

次に，一般の有限次元ベクトル空間で使う形を定める．

% \begin{defi}[有限次元ベクトル空間のテンソル積]
% 複素ベクトル空間$V,W$の基底を，それぞれ
% $\{e_i\}_{i=1}^{m}$，$\{f_j\}_{j=1}^{r}$とする．
% $V\otimes W$を，記号$e_i\otimes f_j$を基底にもつベクトル空間とする．
% また，$v=\sum_i a_ie_i$，$w=\sum_j b_jf_j$に対して
% \begin{equation}
%  v\otimes w
%  :=\sum_{i=1}^{m}\sum_{j=1}^{r}a_ib_j(e_i\otimes f_j)
%  \label{eq:vector-space-tensor-product-definition}
% \end{equation}
% と定める．
% \end{defi}

% この定義は，座標を選んだときの式\eqref{eq:vector-tensor-product-coordinate}と一致する．
特に，テンソル積は
\begin{align}
  (v_1+v_2)\otimes w&=v_1\otimes w+v_2\otimes w,\\
  v\otimes(w_1+w_2)&=v\otimes w_1+v\otimes w_2,\\
  (cv)\otimes w&=v\otimes(cw)=c(v\otimes w)
\end{align}
を満たす．
% $v\otimes w$の形のベクトルを純テンソルという．
% 一般のテンソルは純テンソルの有限個の和で表される．
% ただし，一つの純テンソルで表せるとは限らない．

% \begin{prop}
% ベクトルの族$\{e_i\}_{i=1}^{r}$を$V$の基底，$\{f_j\}_{j=1}^{s}$を$W$の基底とすると，
% $\{e_i\otimes f_j\}_{i,j}$は$V\otimes W$の基底である．
% したがって$\dim(V\otimes W)=rs$である．
% \end{prop}
% \begin{proof}
% 定義より，$e_i\otimes f_j$は$V\otimes W$の基底である．
% その個数は$rs$個であるため，$\dim(V\otimes W)=rs$となる．
% また，任意のベクトル$v=\sum_i a_ie_i$，$w=\sum_jb_jf_j$について
% \[
%   v\otimes w=\sum_{i,j}a_ib_j(e_i\otimes f_j)
% \]
% と表される．これは，式\eqref{eq:vector-space-tensor-product-definition}そのものである．
% \end{proof}

% 内積空間$V,W$に対し，純テンソル上で
% \begin{equation}
%   \langle v_1\otimes w_1,v_2\otimes w_2\rangle
%   :=\langle v_1,v_2\rangle\langle w_1,w_2\rangle
%   \label{eq:tensor-inner-product}
% \end{equation}
% と定め，線形性によって$V\otimes W$全体へ拡張する．
% もとの基底が正規直交基底なら
% \[
%  \langle e_i\otimes f_j,e_k\otimes f_l\rangle
%  =\delta_{ik}\delta_{jl}
% \]
% であるから，積基底も正規直交基底である．

% \begin{defi}[線形写像のテンソル積]
% 線形写像$A:V\to V'$，$B:W\to W'$に対し，
% $(A\otimes B)(v\otimes w)=Av\otimes Bw$を満たす線形写像を
% $A\otimes B$と定める．
% \end{defi}

% \begin{prop}
% 写像の定義域と値域が適合するとき
% \begin{align}
%  (A\otimes B)(C\otimes D)&=AC\otimes BD,\label{eq:tensor-product-composition}\\
%  (A\otimes B)^\dagger&=A^\dagger\otimes B^\dagger
%  \label{eq:tensor-product-adjoint}
% \end{align}
% が成り立つ．したがって，$A,B$がユニタリなら$A\otimes B$もユニタリである．
% これらは，行列のテンソル積の基本的な性質である
% \cite[75頁，定理2.8]{tsukano2023}．
% \end{prop}
% \begin{proof}
% 純テンソル$v\otimes w$に作用させると
% \[
%  (A\otimes B)(C\otimes D)(v\otimes w)
%  =ACv\otimes BDw=(AC\otimes BD)(v\otimes w)
% \]
% である．純テンソルが空間を生成するため式\eqref{eq:tensor-product-composition}が従う．
% また純テンソルについて
% \begin{align*}
%  \langle v'\otimes w',(A\otimes B)(v\otimes w)\rangle
%  &=\langle v',Av\rangle\langle w',Bw\rangle\\
%  &=\langle A^\dagger v',v\rangle\langle B^\dagger w',w\rangle\\
%  &=\langle(A^\dagger\otimes B^\dagger)(v'\otimes w'),v\otimes w\rangle
% \end{align*}
% であり，線形性から式\eqref{eq:tensor-product-adjoint}を得る．
% 写像$A,B$がユニタリなら両式より
% $(A\otimes B)^\dagger(A\otimes B)=I\otimes I$である．
% \end{proof}

複数の量子系を合わせたものを複合量子系という．
複合量子系の状態空間は，各量子系の状態空間のテンソル積で表す．
たとえば2量子ビットの計算基底は
\[
 \ket{00},\ket{01},\ket{10},\ket{11},
 \qquad \ket{ij}:=\ket{i}\otimes\ket{j}
\]
である．本論文では左側の因子を上位の量子ビットとし，
$\ket{x_{n-1}\cdots x_0}$を整数$\sum_{j=0}^{n-1}2^jx_j$に対応させる．

\subsection{複数量子ビットに作用するゲート}

\begin{defi}[制御ユニタリ]
一方の量子ビットを制御量子ビットとする．
また，状態空間$\mathcal K$で表される側を標的とする．このとき，
\begin{equation}
 C(U)=\ket{0}\bra{0}\otimes I_{\mathcal K}
      +\ket{1}\bra{1}\otimes U
 \label{eq:controlled-unitary}
\end{equation}
を制御$U$演算という．
\end{defi}

\begin{prop}
写像$U$がユニタリなら$C(U)$もユニタリである．
\end{prop}
\begin{proof}
射影を$P_0=\ket{0}\bra{0}$，$P_1=\ket{1}\bra{1}$とおくと，
$P_iP_j=\delta_{ij}P_i$，$P_0+P_1=I$である．よって
\begin{align*}
 C(U)^\dagger C(U)
 &=(P_0\otimes I+P_1\otimes U^\dagger)
   (P_0\otimes I+P_1\otimes U)\\
 &=P_0\otimes I+P_1\otimes U^\dagger U
  =(P_0+P_1)\otimes I=I
\end{align*}
である．同様に$C(U)C(U)^\dagger=I$も得られる．
\end{proof}

標的に作用する演算を$U=X$としたものがCNOTゲートである．
また，制御量子ビットが$\alpha\ket{0}+\beta\ket{1}$であり，
標的の状態が$\ket{\psi}$であるとする．線形性より
\begin{equation}
 C(U)(\alpha\ket{0}+\beta\ket{1})\ket{\psi}
 =\alpha\ket{0}\ket{\psi}+\beta\ket{1}U\ket{\psi}
 \label{eq:controlled-unitary-superposition}
\end{equation}
となる．この式は位相推定と制御回転で繰り返し用いる．
